\documentclass[12pt,a4paper]{memoir}

\usepackage[utf8]{inputenc}
\usepackage[T1]{fontenc}
\usepackage{geometry}
\geometry{margin=1in,headheight=14.5pt}

\usepackage{lmodern}
\usepackage{microtype}
\usepackage{mathtools}
\usepackage{amsmath,amssymb,amsthm}
\usepackage{booktabs}
\usepackage{graphicx}
\usepackage[font=small]{subcaption}
\usepackage{tabularx}
\usepackage[svgnames]{xcolor}
\usepackage[hypertexnames=false]{hyperref}
\usepackage[numbers]{natbib}
\bibliographystyle{plain}
\usepackage{cleveref}

% memoir native headers (fancyhdr conflicts with memoir)
\makepagestyle{myheader}
\makeevenhead{myheader}{}{}{}
\makeoddhead{myheader}{}{}{}
\makeevenfoot{myheader}{}{\thepage}{}
\makeoddfoot{myheader}{}{\thepage}{}
\copypagestyle{plain}{myheader}
\pagestyle{myheader}

% memoir abstract: define after \begin{document}

\def\reals{{\mathbb R}}

% memoir chapter style: plain
\chapterstyle{plain}

% ========== DOCUMENT ==========
\begin{document}

% ===== TITLE PAGE =====
\thispagestyle{plain}
\vspace*{2cm}
\begin{center}
    {\LARGE\bfseries Do Standard Neural Networks Learn\\}
    {\LARGE\bfseries Local Block-Spin Renormalization Group Maps?\par}
    \vspace{1.5cm}

    {\Large\bfseries A Controlled Negative Result on the 2D Ising Majority-Vote Task\par}
    \vspace{2cm}

    \begin{tabular}{ll}
        \textbf{Author:} & Xeon Pitar \\
        \textbf{Affiliation:} & Panda Intelligence Software LTD
    \end{tabular}
    \vspace{2cm}

    \parbox{0.9\textwidth}{\small
        \textbf{Abstract:} We investigate whether standard neural networks (MLP with GELU
        activations) exhibit \emph{RG-like advantages} on the 2D Ising
        majority-vote block-spin map, a deliberately narrow local benchmark
        whose target is well approximated by a linear-threshold rule.  The main
        supported evidence in this revision is a unified 10-seed same-scale
        benchmark.  At the critical point $\beta_c$ and lattice size $L=16$,
        the MLP achieves test MSE $0.62 \pm 0.65$, while a single linear layer
        achieves $0.17 \pm 0.09$ on the same task; the linear model's lower mean
        error is clear, but the strength of the difference is test-sensitive
        because the MLP variance is large (Welch $p=0.055$, Mann-Whitney
        $p=0.026$, Cohen $d=-0.98$).  Some off-critical same-scale settings show
        an even stronger linear advantage, e.g.\ $d=-8.36$ at
        $(L,\beta)=(16,0.30)$ where both models have small seed variance.  These
        results support a scoped negative claim: on this local majority-vote
        task, extra fully connected non-linearity does not yield an RG-like
        advantage over a linear baseline.  We additionally ran a true
        whole-lattice tiled transfer experiment under a reduced exploratory
        budget; it is methodologically informative but not decisive, and it does
        not exhibit a stable monotone scale-degradation law.  Dedicated sampler
        diagnostics further show strong critical slowing down, with
        $\tau_{\mathrm{int}}(|M|)\approx 37$ sweeps at $(L,\beta)=(16,\beta_c)$,
        so all cross-scale claims are treated as exploratory rather than
        decisive evidence for RG-like behavior.  A translation-equivariant CNN
        reaches near-zero MSE in most tested settings, but not all, showing that
        locality and translation structure are sufficient to learn the
        supervised block-spin map, but are not by themselves evidence of RG-flow
        learning.  Random-label pilots and a compressed turbulence feasibility
        note are retained only in the appendix.} 
\end{center}
\clearpage
\pagenumbering{roman}
\tableofcontents
\clearpage
\pagenumbering{arabic}

% ===== SECTION 1: INTRODUCTION =====
\section{Introduction}
\label{sec:intro}

The renormalization group (RG), developed by Wilson~\citep{wilson1975}
and crystallized in Kadanoff's block-spin picture~\citep{kadanoff1966},
describes how physical systems behave across scales.  At a critical point,
the RG transformation converges to a fixed point; the approach to this
fixed point is governed by critical exponents~\citep{goldenfeld1992,zinn2021}.

The block-spin picture provides an intuitive mechanical picture of RG.
Imagine a lattice of spins, each pointing either up (+1) or down (-1).  The RG
transformation groups nearby spins into blocks --- for a $2\times2$ block
size $\Lambda$, four adjacent spins become one coarse-grained spin.  To
determine the new spin, we take a majority vote: if three or four spins in the
block point up, the coarse-grained spin points up; if three or four point
down, the coarse-grained spin points down.  This reduces an $L\times L$
lattice to an $L/2\times L/2$ lattice.  Repeating this operation at multiple
scales generates the RG flow.  At each step the system is described by fewer
degrees of freedom but should represent the same physics at longer
wavelengths.

At the \emph{critical temperature} $\beta_c$, something remarkable
happens: fluctuations occur at \emph{all} scales simultaneously.  A spin at
the critical point is correlated with other spins at all distances ---
correlation functions decay as a power law, not exponentially.  There is no
characteristic length scale in the system (the correlation length $\xi$
diverges to infinity).  This is why the critical point is where RG fixed
points live: at $\beta_c$, the system looks statistically the same no matter
what magnification you use.  Zoom in on a critical system and you see the same
pattern of large and small fluctuations you saw before zooming.

The "fixed point" in RG language is a system that is \emph{invariant}
under the coarse-graining transformation.  If you apply the block-spin
transformation to a system already at the fixed point, you get back a system
with the same statistical properties --- same correlation functions up to a
rescaling of length.  The fixed point is the intersection of all RG flows
starting from different microscopic parameters (different temperatures, fields,
couplings).  Physical quantities near the fixed point obey power laws;
the exponents of these power laws are the critical exponents.  They are
universal: they depend only on the dimensionality and symmetry of the order
parameter, not on the microscopic lattice details.

The question of whether neural networks implement RG-like
computations is therefore a question about whether they discover, through
training alone, this hierarchical coarse-graining procedure.  If a standard
architecture trained on physical data implicitly learns to block-average
across scales, this would constitute a strong inductive bias toward the kind
of scale-invariant reasoning that RG formalizes.

The hypothesis that neural networks might \emph{implement} RG-like computations
was proposed by~\citep{mehta2014renormalization}, who observed that restricted
Boltzmann machines learn effective couplings matching RG flow in the 2D Ising
model.  Subsequent work formalized this idea through information-theoretic and
explicit coarse-graining constructions~\citep{kochlug,li2018neuralrg,
pang2020neural,hu2020holographic,bachtis2021}.

We take a different stance: can \emph{standard} architectures (MLP with GELU,
no explicit RG structure) learn block-spin transformations, and do they show
quantifiable signatures of RG-like behavior?  Our emphasis is deliberately
narrow: the target is the local majority-vote block-spin map, not a general
non-linear RG flow.  This question still matters for the interpretability of
learned coarse-graining and surrogate models, but any conclusion must be read
at that same scoped level.

The broader significance extends beyond the Ising model.  If standard neural
networks already implement RG-like computations, this would partially explain
their striking success in modeling multi-scale physical systems such as
turbulent flows~\citep{pope2000,frisch1995}, crack propagation in heterogeneous
materials~\citep{germano1991,smagorinsky1963}, and material defect dynamics.
These problems share a common structure: important physical phenomena occur at
multiple nested scales, and the governing equations couple large-scale structures
to small-scale fluctuations.  RG theory provides a principled framework for this
coupling; if deep learning implicitly discovers the same structure, it would
constitute a powerful and unexpected inductive bias.  Conversely, if standard
architectures do \emph{not} exhibit RG-like behavior, this raises the
fundamental question of what \emph{other} inductive biases they do rely upon ---
and whether those biases are sufficient for robust extrapolation in physical
science.  This connects to the broader research program of "machine learning
meets theoretical physics"~\citep{mehta2014renormalization,pang2020neural,
li2018neuralrg,kochlug}, which seeks to understand the relationship between the
structure of physical law and the structure of learned representations.

\subsection{Research Questions}

\begin{enumerate}
    \item Can a standard MLP learn the block-spin RG transformation
        $T_\Lambda: \{S_i^{(L)}\} \mapsto \{S_i^{(L/\Lambda)}\}$ for the 2D Ising model?

        Here, ``learn'' means \emph{supervised regression}: we generate examples
        of fine-grained spin configurations $S^{\mathrm{fine}}$ along with their
        coarse-grained labels $S^{\mathrm{coarse}}$ computed via the majority-vote
        block-spin transformation (Eq~\ref{eq:block_spin}), and train the network
        to predict $S^{\mathrm{coarse}}$ from $S^{\mathrm{fine}}$.  This is not an
        unsupervised flow-based model nor a generative model; it is a direct
        function approximation task.  The question is whether a standard MLP
        architecture can represent the block-spin map without having it hardwired
        into the architecture.

    \item Is this learning statistically meaningful (vs. random baselines)?

        Any sufficiently expressive model can memorize a finite training set.
        Our original design therefore included a random-label sanity check.
        In the present revision that control remains appendix-level because it
        was produced under a separate 3-seed pilot protocol and is not part of
        the benchmark of record.

    \item Does the learning quality degrade with scale distance, as RG theory predicts?

        RG theory's central quantitative prediction is that information about
        the microscopic scale is progressively lost under repeated coarse-graining.
        If a network is trained on block-spin targets at one lattice size $L$
        and then tested at a different size $L'$, the test error should increase
        monotonically with $|L - L'|$, because the RG flow has moved further
        away from the training scale.  This cross-scale transfer is \emph{the}
        defining operational test of RG-like behavior in our framework.  It is
        qualitatively distinct from within-scale learning (which merely confirms
        the network can represent the block-spin map at a single scale) and
        requires the network to have discovered the \emph{hierarchical}
        structure of scale-coupling, not just a local majority-vote rule.  In
        the present revision we therefore treat monotone degradation as a
        hypothesis whose operationalization remains unstable for the current
        local-rule benchmark, rather than as a settled diagnostic.

    \item Does the framework extend to a neural turbulence closure reproducing
        Kolmogorov spectrum scaling?

        Kolmogorov's 1941 hypothesis~\citep{kolmogorov1941} for fully developed
        turbulence states that the energy spectrum in the inertial range follows
        $E(k) \propto k^{-5/3}$, where $k$ is the wavenumber.  This power law
        emerges from a cascade: large eddies inject energy at the forcing scale,
        break into smaller eddies, and so on until viscosity dissipates at the
        Kolmogorov microscale.  The $-5/3$ exponent is remarkable because it can
        be derived almost entirely from dimensional analysis plus the assumption
        of universality (the inertial range statistics depend only on the energy
        cascade rate $\epsilon$).  If a neural network trained \emph{only} on
        subgrid stress closures (the small-scale fluctuations unresolved by the
        simulation grid) can reproduce this exponent, it would constitute strong
        evidence that the network has learned an RG-like scale-coupling in a
        physically non-trivial setting beyond the controlled Ising testbed.
        This would be striking because the network never sees the inertial range
        directly during training; the correct exponent would have to emerge from
        the physics embedded in the training data.
\end{enumerate}

\subsection{Operational Targets and Hypotheses}
\label{sec:operational_def}

We retain the labels R1--R3 to organize the paper, but in this revision they
do not all carry the same evidentiary status.  R1 motivates the original
random-label sanity check, R2 is best read as a \emph{transfer hypothesis} to
be tested under a unified multi-scale protocol, and R3 is an exploratory
physics-extension target.

\medskip\noindent
Concretely:
\begin{enumerate}
    \item[(R1)] \textbf{Appendix-level sanity check:} Test MSE on block-spin prediction
        is significantly better than a random-label baseline at the same $\beta$.
    \item[(R2)] \textbf{Scale-transfer hypothesis:} Under a unified multi-scale
        transfer protocol, a genuinely RG-like model should exhibit a stable
        relationship between transfer quality and scale distance and should
        outperform competitive linear baselines.
    \item[(R3)] \textbf{Turbulence extension target:} A neural closure model trained
        on subgrid stress reproduces Kolmogorov $k^{-5/3}$ inertial-range scaling
        in the energy spectrum.
\end{enumerate}

A convincing positive claim of RG-like computation would require aligned
evidence across these targets under matched protocol and budget.  The present
revision provides strong same-scale evidence, but only exploratory evidence for
R2 and R3.

\textbf{Revision scope.}  The present revision separates results by protocol.
The primary quantitative claims use only the unified 10-seed benchmark saved in
\path{outputs/rg_bench/cross_scale/}.  An earlier fixed-split pilot
ablation (including random-label controls and width sweeps) used a different
3-seed protocol; it is retained only as a sanity check and is \emph{not}
numerically mixed with the primary benchmark tables.  A separate
whole-lattice tiled transfer experiment and dedicated sampler diagnostics are
reported with their own reduced exploratory budget and interpreted
accordingly.

\textbf{What this paper establishes, and what it does not:}
This paper locates itself as a scoped negative-result paper.
The strongest supported claim is the same-scale comparison at
$L=16$, $\beta_c$: the linear baseline outperforms the MLP
in mean test MSE ($0.17 \pm 0.09$ vs.\ $0.62 \pm 0.65$), although the
strength of this difference is test-sensitive because of the MLP's large
variance.  The true whole-lattice tiled
transfer diagnostic does not support the original monotone-degradation R2
hypothesis, but the current reduced-budget results are better interpreted as a
warning that R2 is not yet a stable operational discriminator in this local-rule
setting.  Combined with the observed sampler autocorrelation at criticality, we
treat cross-scale conclusions as exploratory and unresolved.  We do not make a
primary quantitative R1 claim from the separate random-label pilot.  R3
(turbulence) remains a feasibility demonstration only and does not count as
evidence for any criterion.

The main negative finding is: a linear model substantially outperforms
the MLP on the same block-spin task, consistent with the majority-vote
transformation being dominated by a local linear-threshold rule.  This is a precise
negative result scoped specifically to \emph{locally linear block-spin
tasks}.  It does not imply that standard deep architectures cannot
learn RG-like computations in general; that would require testing on
tasks where the RG transformation is genuinely non-linear in the
couplings (e.g., non-zero external field, frustrated magnets, or quantum
systems).  It also does not contradict the RBM-based positive findings
of Mehta \& Schwab (2014), which rely on an architectural inductive bias
--- the bipartite energy function $E(v,h) = -b^\top v - c^\top h - v^\top W h$ ---
that is absent from standard MLPs.  The scope of our claim is limited:
standard MLPs do not exhibit RG-like advantages on majority-vote block-spin
tasks that are well-approximated by linear projections.

% ===== SECTION 2: BACKGROUND ======
\section{Background: Ising Model and Block-Spin RG}
\label{sec:bg}

The 2D Ising model on an $L \times L$ square lattice with periodic
boundary conditions has Hamiltonian:
\begin{equation}
    H = -J\sum_{\langle ij\rangle} S_i S_j - h\sum_i S_i,
    \qquad S_i \in \{-1, +1\}.
    \label{eq:ising}
\end{equation}
We set $J=1$, $h=0$ (zero field).  The critical inverse temperature is
Onsager's exact value~\citep{onsager1944}:
\begin{equation}
    \beta_c = \frac{1}{2}\log(1+\sqrt{2}) \approx 0.4407.
    \label{eq:onsager}
\end{equation}

The 2D Ising model holds a special place in the RG literature for three
interrelated reasons.  First, Lars Onsager's 1944 exact solution is one of the
very few non-trivial phase transitions in statistical mechanics that can be
solved analytically in closed form.  This provides \emph{ground truth} for all
relevant quantities: the free energy, correlation functions, and critically,
the critical exponents $\nu=1$, $\eta=1/4$, $\beta=1/8$.  Any numerical method
for studying RG must be tested against these exact values.  Second, the 2D
Ising model is the \emph{simplest} model with a non-trivial continuous phase
transition: it has a paramagnetic phase at high temperature, a ferromagnetic
phase at low temperature, and a second-order transition at $\beta_c$ with
well-defined critical exponents.  It captures the essential physics of the
ferromagnetic-paramagnetic transition without the complications of additional
order parameters, anisotropy, or frustration.  Third, the 2D Ising universality
class is remarkably robust: any spin model with short-range interactions and a
discrete symmetry ($Z_2$ inversion) on a 2D lattice belongs to the same
universality class and shares the same critical exponents.  This universality
means the 2D Ising results generalize broadly, making it an ideal testbed for
controlled theoretical experiments~\citep{goldenfeld1992,amit2005,cardy1996}.

The block-spin RG transformation with block size $\Lambda=2$ is:
\begin{equation}
    S'_{\lfloor i/\Lambda \rfloor, \lfloor j/\Lambda \rfloor}
    = \mathrm{sgn}\!\Big(\sum_{k,l \in \mathrm{block}(i,j)} S_{k,l}\Big),
    \label{eq:block_spin}
\end{equation}
i.e., majority vote within each $2\times2$ block.  Iterating
Eq~\eqref{eq:block_spin} generates an RG flow.  At $\beta_c$ the flow
converges to the Wilson-Fisher fixed point; away from $\beta_c$ it flows
to ordered or disordered phases.

The majority-vote construction deserves careful explanation.  The Ising order
parameter is the \emph{magnetisation} $m = \langle S_i \rangle$, a sum over spins
that is discontinuous (non-analytic) at $\beta_c$ in the thermodynamic limit.
At the critical point, the block-averaged spin $\sum_{k,l \in \mathrm{block}} S_{k,l}$
is a continuous random variable (it takes values in $\{-4,-2,0,2,4\}$), but the
magnetisation of the coarse-grained system is the sign of this sum.  The sign
function is a non-linear threshold that implements the discontinuous phase
transition: below $\beta_c$ the sum is overwhelmingly positive (ordered phase,
all spins aligned), while above $\beta_c$ the sum fluctuates around zero
(disordered phase).  At $\beta_c$ itself, the sum can be positive or negative
with roughly equal probability, and the sign function determines whether the
coarse-grained block falls in the up or down domain.  Geometrically, the sign
function partitions the 4-dimensional spin space ($\pm1$ per spin) into two
hemispheres separated by the plane $\sum_i S_i = 0$.  This is a \emph{linear}
threshold in the space of block-summed features, which explains why the
majority-vote transformation is well-approximated by a linear model ---
the non-linearity enters only through the sign function applied to a single
linear combination of the four spins in each block.

Critical exponents for the 2D Ising universality class are exactly known:
$\nu=1$, $\eta=1/4$, $\beta=1/8$~\citep{onsager1944,kadanoff1967,stanley1971}.

Each critical exponent governs a specific physical quantity near $\beta_c$.
The exponent $\nu$ controls how fast the correlation length $\xi$ diverges as
the critical point is approached from either side: $\xi \sim |\beta - \beta_c|^{-\nu}$.
Physically, near $\beta_c$ fluctuations become large and long-range; $\nu=1$
in 2D Ising means this divergence is relatively slow.  The exponent $\eta$
characterises the decay of the spin-spin correlation function at the critical
point itself: $G(r) \sim r^{-(d-2+\eta)}$ with $\eta=1/4$.  In 2D this gives
$G(r) \sim r^{-1.75}$, a power law slower than the mean-field prediction
$r^{-1}$ because strong fluctuations in 2D modify the correlation structure.
The exponent $\beta$ governs the rise of the order parameter (magnetisation)
below $\beta_c$: $m \sim (-\epsilon)^{\beta}$ where $\epsilon = (\beta_c - \beta)/\beta_c$
is the reduced temperature.  The value $\beta=1/8$ (much smaller than the
mean-field value $\beta=1/2$) reflects the difficulty of ordering a 2D system
against thermal fluctuations.  These three exponents are not independent;
they satisfy Rushbrooke's relation $2-\alpha = d\nu$, Josephson scaling
$\nu d = 2-\eta + \beta(\delta-1)$ and others, forming a tightly constrained
set that the Onsager solution uniquely determines~\citep{goldenfeld1992,amit2005}.

% ===== SECTION 3: FRAMEWORK =====
\section{Neural Block-Spin Learning Framework}
\label{sec:framework}

We frame block-spin learning as a supervised regression problem:
given spin configurations $S^{\mathrm{fine}} \in \{-1,+1\}^{L\times L}$,
predict $S^{\mathrm{coarse}} \in \{-1,+1\}^{L/2\times L/2}$ from majority-vote
block-spin RG (Eq~\eqref{eq:block_spin}).

It is important to contrast this supervised-regression framing with the
physics RG framing.  In theoretical physics, RG is a \emph{functional
renormalization group flow} --- a continuous transformation of the coupling
constants $K$ in the Hamiltonian under change of cutoff scale.  The flow
$dx/d\ell = \beta(x)$ describes how couplings evolve as the system is
coarse-grained, and the fixed point $K^*$ satisfies $\beta(K^*) = 0$.  This
is a continuous, deterministic dynamical system on the space of coupling
constants~\citep{wilson1974renormalization,wilson1975,goldenfeld1992,zinn2021}.
Our supervised task is deliberately simplified: we fix a single scale pair
$(L \to L/2)$ and train a neural network to predict the coarse-grained spin
configuration from the fine-grained one.  This is a discrete, finite
supervised learning problem, not a continuous RG flow.  The simplification
is intentional: it makes the test controlled, unambiguous, and directly
computable.  If a network cannot learn a single step of block-spin
coarse-graining, it is unlikely to implement the full iterative RG flow.
Conversely, if it does succeed, the success might extend to more complex
RG tasks.  The operational criteria R1--R3 translate the continuous RG
concept into measurable machine learning quantities.

\subsection{Network Architecture}

The MLP architecture (3 hidden layers, GELU activations, $\tanh$ output for
spin-like bounded output) is deliberately standard:
\begin{equation}
    \text{MLP}: \reals^{L^2} \xrightarrow{\text{GELU}} \reals^{256}
    \xrightarrow{\text{GELU}} \reals^{256}
    \xrightarrow{\text{GELU}} \reals^{128}
    \xrightarrow{\text{GELU}} \reals^{64}
    \xrightarrow{\tanh} \reals^{(L/2)^2}.
    \label{eq:arch}
\end{equation}
The encoder (first three layers) compress $L^2$ inputs; the decoder
produces $(L/2)^2$ coarse-grained spins.  No explicit RG structure is
hardwired.  Training: Adam optimizer, lr=$10^{-3}$, MSE loss,
batch size 32, 200 epochs.  The unified benchmark reported in the revised
tables uses seeds $\{42,\dots,51\}$ ($n=10$).  A separate fixed-split pilot
ablation uses 3 seeds and is identified explicitly when mentioned.  Statistical
inference for the primary benchmark uses the saved summaries in
\texttt{cross\_scale\_statistics.json}: bootstrap 95\% CIs for per-model mean
MSE, plus Welch tests and Cohen's $d$ for MLP-vs-Linear comparisons at fixed
$(L,\beta)$.

The choice of this specific MLP architecture deserves justification.  We
deliberately chose a \emph{standard} feedforward architecture with no
explicit RG structure --- no weight-sharing across scales, no hierarchical
inductive bias, no symmetry constraints.  The GELU activation~\citep{hendrycks2016gaussian}
is a smooth approximation to the ReLU that is standard in modern deep learning.
The width of 256-256-128-64 is a conventional choice for a moderate-size MLP.
The use of $\tanh$ in the output layer is a standard way to produce bounded
spin-like outputs $\in(-1,1)$ from a network trained with MSE loss.  The
central question is whether \emph{any} of these standard architectural
choices implicitly bias the network toward RG-like computations.  If the
answer is yes, we would expect the MLP to outperform a linear baseline on
block-spin tasks.  If the answer is no, then RG-like inductive biases must
be explicitly engineered (as in Neural RG papers), not discovered through
standard architecture design.

The linear baseline serves as the minimal model for this task.  A single
linear layer from $\reals^{L^2}$ to $\reals^{(L/2)^2}$ has exactly $(L^2)(L/2)^2$
parameters (e.g., for $L=16$: $256 \times 64 = 16{,}384$ weights).  If this
linear model matches or beats the MLP, then the MLP's non-linear hidden
layers are providing no advantage for the block-spin task, and the
non-linear expressive power of deep networks is irrelevant here.  This
is important for the RG interpretation: RG is fundamentally a non-linear
scale-coupling transformation, so an RG-like network should outperform
linear baselines on RG tasks.  Finding the opposite confirms that the
block-spin transformation at $L=16$ is approximately linear in the spin
variables.

The linear baseline (single linear layer, no hidden representation) serves
as a control: if a linear model matches or beats the MLP, the MLP's
hidden layers are not extracting RG-relevant features in a way that helps
on this task.

\subsection{Heuristic: Majority Vote as a Local Linear-Threshold Map}
\label{sec:linear-rg-approx}

Let $T : \{-1,+1\}^{L\times L} \to \{-1,+1\}^{L/2\times L/2}$
denote the exact block-spin RG transformation:
\begin{equation}
  T(s)_{IJ} \;=\; \operatorname{sign}\!\Bigl(\sum_{i\in B_{IJ}} s_i\Bigr),
  \qquad B_{IJ} = \{(2I\!-\!1,\!2J\!-\!1),\ldots,(2I,\!2J)\},
\end{equation}
and let
\begin{equation}
  z_{IJ}(s) \;=\; \sum_{i\in B_{IJ}} s_i \in \{-4,-2,0,2,4\}
\end{equation}
be the block sum.  Each coarse spin is therefore produced by applying a
threshold to a \emph{linear} score on a local $2\times2$ block:
\begin{equation}
  T(s)_{IJ} = \operatorname{sign}\!\bigl(z_{IJ}(s)\bigr).
\end{equation}

This observation is only a heuristic.  It does \emph{not} imply a uniform
approximation theorem in $L$ or $\beta$, and we make no such claim in this
revision.  The point is more modest: a shallow model can compute each block sum
exactly with a sparse linear map, leaving only the final threshold
non-linearity.  That local linear-threshold structure makes it plausible that a
linear baseline can be competitive on the supervised block-spin task.
Whether that heuristic predicts actual test-time performance remains an
empirical question, answered by the benchmark in Section~\ref{sec:ising_exp}.

% ===== SECTION 4: ISING EXPERIMENTS =====
\section{Experiments: 2D Ising Model}
\label{sec:ising_exp}

All experiments use single-spin Metropolis-Hastings Monte Carlo with
$N_{\mathrm{train}}=500$, $N_{\mathrm{test}}=300$, equilibration $10^3$
steps~\citep{metropolis1953,hastings1970}.

Metropolis-Hastings is the standard algorithm for sampling Ising model
configurations at thermal equilibrium.  Starting from an initial spin
configuration, the algorithm proposes flipping a single spin, accepts the flip
with probability $\min(1, \exp(-\Delta E/T))$ where
$\Delta E$ is the energy change, and rejects it otherwise.  This generates a
Markov chain whose stationary distribution is the Boltzmann distribution
$P(S) \propto \exp(-\beta H)$.  The equilibration period of $10^3$ steps
is a pragmatic burn-in choice rather than a proof of mixing.  In the current
scripts we record successive post-burn-in states without an explicit
integrated-autocorrelation-time estimate or thinning interval.  At $\beta_c$,
critical slowing down implies that within-chain samples remain correlated, so
we do \emph{not} claim that the resulting configurations are independent draws
from the thermodynamic ensemble.  Separate train/test chains reduce direct data
reuse between splits, but the reported error bars should be interpreted as
benchmark variability under this sampler budget, not as a certified
independent-sample uncertainty quantification.  A dedicated diagnostic run
reported in Appendix~\ref{app:mixing} gives
$\tau_{\mathrm{int}}(E)=6.9 \pm 1.9$ and
$\tau_{\mathrm{int}}(|M|)=9.6 \pm 3.0$ sweeps at $(L,\beta)=(8,\beta_c)$, and
$18.8 \pm 7.3$ and $37.0 \pm 15.1$ sweeps at $(16,\beta_c)$, whereas away
from criticality the corresponding values remain near $2$--$3$ sweeps.  This
quantifies the critical slowing-down caveat rather than leaving it qualitative.
Cluster algorithms such as Swendsen-Wang and Wolff would generally be better
suited to this regime, but they are not used in the current benchmark
implementation~\citep{swendsen1987,wolff1989,sokal1997,newman1999,binder2010}.
The appropriate interpretation of the reported error bars and significance tests
is therefore that of a finite-budget benchmark under correlated sampling,
rather than thermodynamic-ensemble-level inference under effectively
independent draws.  At $\beta_c$, the corresponding effective sample sizes are
only approximately $218$ and $156$ for $(L,\beta)=(8,\beta_c)$, and
$80$ and $41$ for $(16,\beta_c)$, for $E$ and $|M|$ respectively.

\paragraph{Evidence tiers.}
To keep evidence levels explicit, Table~\ref{tab:evidence_tiers} summarizes the
three protocol tiers used in this revision.

\begin{table}[ht]
    \centering
    \caption{Evidence tiers used in the revised manuscript.}
    \small
    \begin{tabular}{@{}p{0.20\textwidth} p{0.28\textwidth} p{0.22\textwidth} p{0.22\textwidth}@{}}
        \toprule
        \textbf{Tier} & \textbf{Protocol} & \textbf{Budget} & \textbf{Role} \\
        \midrule
        Main benchmark & Same-scale Ising benchmark & $n=10$, $500/300$, 200 epochs & Primary evidence \\
        Exploratory transfer & Whole-lattice tiled transfer & $n=3$, $200/100$, 40 epochs & Methodological clarification \\
        Appendix diagnostics & Pilot / turbulence / mixing & mixed auxiliary budgets & Non-primary support \\
        \bottomrule
    \end{tabular}
    \label{tab:evidence_tiers}
\end{table}

\subsection{Result 1: Same-Scale Block-Spin Benchmark}
\label{sec:block_spin_learning}

\textbf{Setup.}  The primary quantitative claims in this revision use the
unified 10-seed benchmark saved in
\texttt{outputs/rg\_bench/cross\_scale/}.  For each
$L \in \{4,8,16\}$, $\beta \in \{0.30,\beta_c,0.60\}$, and model in
\{MLP, Linear, CNN, RGMLP\}, we train with the same optimizer and data
budget and report held-out test MSE.  Earlier fixed-split pilot ablations
(random-label controls and width sweeps) were generated by a different
3-seed protocol and are therefore \emph{not} mixed numerically with the
tables below.

\textbf{Results.}  Table~\ref{tab:ising_main} and Figure~\ref{fig:block_spin_models}.

\begin{table}[ht]
    \centering
    \caption{Main same-scale Ising benchmark ($n=10$ seeds).
        All models are trained on block-spin labels under the unified protocol
        saved in \texttt{cross\_scale\_statistics.json}.  Values are mean
        $\pm$ std dev over seeds $\{42,\dots,51\}$.  This is the main benchmark
        of record and is not directly comparable in evidence strength to the
        reduced-budget exploratory transfer table below.}
    \resizebox{\textwidth}{!}{%
    \begin{tabular}{lcccc}
        \toprule
        \textbf{Model} & \textbf{$\beta$} & \textbf{$L=4$ MSE} & \textbf{$L=8$ MSE} & \textbf{$L=16$ MSE (same-$L$)} \\
        \midrule
        MLP     & 0.30  & $0.0005 \pm 0.0006$  & $0.1198 \pm 0.0128$  & $0.6006 \pm 0.0287$  \\
        MLP     & $\beta_c$ & $0.0307 \pm 0.0239$  & $0.2017 \pm 0.2187$  & $0.6249 \pm 0.6456$  \\
        MLP     & 0.60  & $0.0168 \pm 0.0231$  & $0.1973 \pm 0.5341$  & $0.1178 \pm 0.1695$  \\[4pt]
        Linear  & 0.30  & $0.0314 \pm 0.0059$  & $0.0536 \pm 0.0051$  & $0.4160 \pm 0.0123$  \\
        Linear  & $\beta_c$ & $0.0357 \pm 0.0265$  & $0.0692 \pm 0.0309$  & $0.1736 \pm 0.0862$  \\
        Linear  & 0.60  & $0.0159 \pm 0.0216$  & $0.0140 \pm 0.0092$  & $0.0106 \pm 0.0099$  \\[4pt]
        CNN     & 0.30  & $3.2\times10^{-5} \pm 0.9\times10^{-5}$ & $1.0\times10^{-5} \pm 0.1\times10^{-5}$ & $3.6\times10^{-6} \pm 0.5\times10^{-6}$ \\
        CNN     & $\beta_c$ & $5.4\times10^{-4} \pm 1.1\times10^{-3}$ & $2.8\times10^{-5} \pm 2.9\times10^{-5}$ & $6.0\times10^{-6} \pm 5.0\times10^{-6}$ \\
        CNN     & 0.60  & $0.0248 \pm 0.0432$  & $0.0451 \pm 0.0714$  & $0.0028 \pm 0.0073$  \\[4pt]
        RGMLP   & 0.30  & $0.4928 \pm 0.3231$  & $0.7612 \pm 0.3163$  & $1.1784 \pm 0.2744$  \\
        RGMLP   & $\beta_c$ & $0.1542 \pm 0.0672$  & $0.5402 \pm 0.3147$  & $1.0646 \pm 1.0284$  \\
        RGMLP   & 0.60  & $0.6601 \pm 1.3218$  & $1.7935 \pm 1.4374$  & $0.9792 \pm 1.2162$  \\
        \bottomrule
    \end{tabular}}
    \label{tab:ising_main}
\end{table}

\begin{figure}[!htbp]
    \centering
    \includegraphics[width=0.82\textwidth]{outputs/rg_bench/cross_scale/figures/model_comparison_bars.png}
    \caption{Within-scale Ising block-spin performance across $\beta$.
        The CNN reaches near-zero MSE in most tested settings, with visible
        degradation only in the low-$L$ ordered cases at $\beta=0.60$, while
        the linear model matches or outperforms the MLP at the critical point.}
    \label{fig:block_spin_models}
\end{figure}

Key findings from Table~\ref{tab:ising_main}:

\begin{enumerate}
    \item[(i)] \textbf{At $L=16$ and $\beta_c$, the linear baseline beats the MLP.}
        The unified same-scale benchmark gives
        $0.1736 \pm 0.0862$ for Linear vs.\ $0.6249 \pm 0.6456$ for MLP at
        $(L,\beta)=(16,\beta_c)$.  The mean ordering favors the linear model,
        but the strength of the difference depends on the test because the MLP
        variance is large (Welch $p=0.055$, Mann-Whitney $p=0.026$,
        Cohen's $d=-0.98$).  The critical-point task is therefore learnable,
        but the extra non-linear capacity of the MLP does not provide a robust
        same-scale advantage.

    \item[(ii)] \textbf{Error depends strongly on phase and lattice size.}
        At $L=16$, the linear model is easiest in the ordered phase
        ($0.0106 \pm 0.0099$ at $\beta=0.60$), while the MLP remains high in
        both the disordered and critical regimes
        ($0.6006 \pm 0.0287$ at $\beta=0.30$ and
        $0.6249 \pm 0.6456$ at $\beta_c$).  The temperature dependence in
        Figure~\ref{fig:temperature_dependence}
        shows that the critical regime carries the largest seed-to-seed
        instability.  In particular, the disordered same-scale setting
        $(L,\beta)=(16,0.30)$ yields an extreme effect size
        ($d=-8.36$) because the mean gap is large while both models have small
        standard deviation.

    \item[(iii)] \textbf{Architectural bias matters more than generic depth.}
        The CNN with explicit locality and translation structure reaches
        near-zero MSE across nearly all $(L,\beta)$ settings, whereas RGMLP is
        unstable and often worse than the plain MLP.  The task is therefore
        perfectly learnable with an appropriate inductive bias, but standard
        fully connected depth alone is not enough.  This is not, by itself,
        evidence that the CNN has learned a physically meaningful RG flow; it
        shows only that locality and translation structure are sufficient for
        the supervised block-spin map.
\end{enumerate}

The earlier random-label control is retained only as an appendix-level pilot
sanity check and is not used for the main quantitative claims in this revision.

Figure~\ref{fig:temperature_dependence} summarizes the within-scale
temperature trend at $L=16$: both models improve away from criticality,
and the ordered phase yields the lowest MSE for both models.

\begin{figure}[!htbp]
    \centering
    \includegraphics[width=0.72\textwidth]{outputs/rg_bench/cross_scale/figures/temperature_dependence.png}
    \caption{Temperature dependence of within-scale block-spin MSE at $L=16$.
        Both models improve away from criticality, and the ordered phase yields
        the lowest MSE for both models.}
    \label{fig:temperature_dependence}
\end{figure}

\clearpage
\subsection{Result 2: Whole-Lattice Tiled Transfer (Exploratory)}
\label{sec:cross_scale}

\textbf{Setup.}  To replace the earlier patch-only proxy with a true
whole-lattice evaluation, we trained source-scale predictors on
$L \in \{4,8\}$ and then tiled the learned local model over the full target
lattice.  Concretely, a model trained on $L=8$ receives each non-overlapping
$8\times8$ patch of a $16\times16$ configuration, predicts the corresponding
$4\times4$ coarse block, and the four predicted blocks are assembled into the
full $8\times8$ coarse lattice.  The same construction yields $4\to8$ and
$4\to16$ transfer.  Because this experiment was added after the main 10-seed
benchmark, we ran it under a reduced exploratory budget
($n=3$ seeds, $N_{\mathrm{train}}=200$, $N_{\mathrm{test}}=100$, 40 epochs)
and compare transfer against same-budget same-$L$ baselines from the same
script.

\begin{table}[ht]
    \centering
    \caption{Exploratory whole-lattice tiled transfer at $\beta_c$.
        Same-budget baselines and transfers come from the
        \texttt{whole\_lattice\_transfer/summary.json}
        using $n=3$ seeds, $N_{\mathrm{train}}=200$, $N_{\mathrm{test}}=100$,
        and 40 epochs.  ``Deg.'' is relative to the same-budget same-$L=16$
        baseline for the same model.}
    \begin{tabular}{lcccc}
        \toprule
        \textbf{Source $\to$ 16} & \textbf{MLP} & \textbf{MLP Deg.} & \textbf{Linear} & \textbf{Linear Deg.} \\
        \midrule
        $16 \to 16$ & $0.463 \pm 0.174$ & $1.00\times$ & $0.287 \pm 0.170$ & $1.00\times$ \\
        $8 \to 16$  & $0.293 \pm 0.190$ & $0.63\times$ & $0.244 \pm 0.129$ & $0.85\times$ \\
        $4 \to 16$  & $0.141 \pm 0.107$ & $0.31\times$ & $0.222 \pm 0.113$ & $0.77\times$ \\
        \bottomrule
    \end{tabular}
    \label{tab:whole_lattice_transfer}
\end{table}

\begin{itemize}
    \item \textbf{The true tiled experiment does not support monotone degradation.}
        At $\beta_c$, neither model deteriorates when transferred from
        $L=8$ to $L=16$ under this tiled evaluation, and the $4\to16$
        transfer is numerically \emph{better} still.  The originally proposed
        monotone scale-distance interpretation therefore fails in this local-rule
        setting.
    \item \textbf{Linear baselines remain fully competitive.}
        Although the MLP improves substantially relative to its same-budget
        $16\to16$ baseline, the linear model also transfers competitively:
        $0.244 \pm 0.129$ for $8\to16$ and $0.222 \pm 0.113$ for $4\to16$,
        versus $0.287 \pm 0.170$ at $16\to16$.
    \item \textbf{The cross-scale story is protocol-sensitive.}
        The new whole-lattice tiled evaluation gives a different qualitative
        picture from the earlier patch-only proxy, which means cross-scale
        conclusions depend strongly on how transfer is operationalized.
        Appendix~\ref{app:whole_lattice} (Table~\ref{tab:whole_lattice_all_beta})
        shows that this non-monotone behavior is not confined to $\beta_c$.
\end{itemize}

\textbf{R2 assessment.}  The new whole-lattice experiment removes the
``top-left patch'' ambiguity of the earlier proxy, but it still does not
deliver a stable positive RG signature.  Instead, it shows that transfer in
this benchmark is highly protocol- and budget-sensitive: the tiled
whole-lattice evaluation does \emph{not} display monotone degradation with
scale distance, and linear baselines remain competitive throughout.  The right
interpretation is therefore not that R2 has been cleanly falsified, but that
for the present local-rule setting the original R2 formulation is not yet a
stable operational discriminator.  We therefore keep R2 unresolved.

\paragraph{Auxiliary analyses moved to the appendix.}
Legacy pilot sanity checks, sampler diagnostics, and the turbulence
feasibility experiment are reported in
Appendices~\ref{app:mixing}--\ref{app:turbulence} and are not used as primary
main-text evidence.

% ===== SECTION 6: DISCUSSION =====
\section{Discussion}
\label{sec:discussion}

\subsection{Why Does the Linear Model Outperform the MLP?}

The strongest supported result in the revised manuscript is the unified
same-scale benchmark at $(L,\beta)=(16,\beta_c)$:
Linear reaches $0.1736 \pm 0.0862$ while MLP reaches
$0.6249 \pm 0.6456$ ($n=10$).  We focus on this result
because it comes from the same 10-seed protocol used throughout the main
tables; the earlier fixed-split random-label pilot is not used for the
primary quantitative claim.

We offer three non-mutually-exclusive interpretations:

\begin{enumerate}
    \item[(i)] \textbf{Target simplicity near $\beta_c$:} The block-spin
        majority-vote at criticality is effectively a \emph{linear}
        threshold function: each coarse spin is the sign of a local sum of
        four $\pm1$ inputs.  A shallow linear map can compute those block sums
        exactly, leaving only the final threshold non-linearity.  This local
        linear-threshold structure makes it plausible that extra hidden-layer
        non-linearities are unnecessary for the supervised task studied here.

    \item[(ii)] \textbf{Feature redundancy:} The first linear layer can
        already compute $2\times2$ block sums, which are sufficient for
        the majority-vote target.  Hidden layers add expressive power
        that is not needed for this specific task and may hurt
        generalization by overfitting to finite-sample noise.  In our setting
        there is no evidence that generic depth discovers a better scale-coupled
        representation than the linear baseline; the CNN result suggests that
        architectural bias, not depth alone, is the missing ingredient.

    \item[(iii)] \textbf{Training dynamics:} The MLP may converge to
        different local minima across seeds, leading to high variance.
        The unified benchmark records standard deviations of $0.646$ for the
        MLP and $0.086$ for the linear model at $(L,\beta)=(16,\beta_c)$.
        This instability is itself part of the negative result: the MLP's
        added flexibility manifests as optimization variance rather than as a
        reliable accuracy advantage.
\end{enumerate}

This finding suggests that, on the present benchmark, the RG-relevant
computation is dominated by a local linear-threshold rule rather than by a
hierarchical non-linear representation.

\subsection{Implications for RG-like Behavior Assessment}

The whole-lattice tiled experiment sharpens the methodological lesson.
Once transfer is evaluated on the full target lattice rather than on a
single patch, the original monotone-degradation criterion still does not
emerge.  At $\beta_c$, both $8\to16$ and $4\to16$ transfer remain
competitive with the same-budget same-$L=16$ baseline, and the smaller source
scale can even perform better numerically.  This means that, in the present
local-rule setting, R2 is not a stable discriminator of RG-like behavior.

Three non-exclusive mechanisms may explain why smaller source scales can appear
to transfer better in the tiled experiment.  First, \emph{local-rule tileability}:
once the source model has learned a sufficiently accurate local majority-vote
rule, that rule can be repeated across the larger lattice without needing a new
global representation.  Second, \emph{source-task simplification}: the
$4\times4$ and $8\times8$ source problems may be easier optimization problems
than the $16\times16$ one, so a smaller source scale can yield a cleaner local
predictor.  Third, \emph{budget mismatch}: under the reduced exploratory budget,
the same-$L=16$ baselines may remain under-optimized relative to the tiled
transfer models, which weakens any attempt to read the ratios as intrinsic RG
flow behavior.  This suggests that tiled transfer may primarily probe reusable
locality of the learned block rule, rather than hierarchical scale-coupling in
the RG sense.

The across-temperature whole-lattice results strengthen this interpretation:
non-monotone transfer is not a critical-point-specific anomaly, but looks more
like a protocol artifact of tiled local-rule reuse across source scales.

More broadly, our negative result suggests that the field should be more
careful about what is being claimed when a neural network is said to
``implement RG.''  If a task is simple enough to be solved by a linear
baseline, then apparent RG-like structure in the learned representation may
be coincidental rather than causal.  Stronger claims should be reserved for
settings in which linear baselines fail, transfer protocols match the
stated operational criterion, and sampling uncertainty is characterized more
carefully than in the present benchmark.

% ===== SECTION 7: RELATED WORK =====
\section{Related Work}
\label{sec:related}

There are now two largely separate literatures around machine learning and
statistical mechanics.  One uses supervised or unsupervised learning to detect
phases, order parameters, and transition points in classical or quantum lattice
systems~\citep{wang2016,carrasquilla2017,vannieuwenburg2017,hu2017,
wetzel2017,wangzhai2017,beach2018,zhang2019percolation,
rodrigueznieva2019,dcn2018}.  More interpretable variants ask what features the
model has actually learned, for example through kernel methods, explicit order
parameters, or symmetry diagnostics~\citep{ponte2017,wetzelscherzer2017,
greitemann2019,liu2019interpretable,decelle2019}.  Our benchmark is narrower
than that literature: we do not ask whether a network can detect a phase
transition, but whether it gains a specifically RG-like advantage on a known
coarse-graining map.

A second literature aims more directly at RG itself.  Mehta and
Schwab~\citep{mehta2014renormalization} connected RBMs with variational RG;
Koch-Janusz and Ringel~\citep{kochlug} recast coarse graining in terms of
mutual information; Li and Wang~\citep{li2018neuralrg}, Pang et
al.~\citep{pang2020neural}, Hu et al.~\citep{hu2020holographic}, and Bachtis
et al.~\citep{bachtis2021} proposed architectures or objectives that encode
multi-scale structure explicitly.  Those papers are valuable precisely because
their model classes are RG-motivated.  Our question is stricter and more
negative-result oriented: if one removes the explicit RG inductive bias and
compares against strong linear baselines, does a standard fully connected MLP
still exhibit an RG-like advantage on the supervised block-spin task?

Finally, our treatment of sampling is informed by the Monte Carlo literature.
Local Metropolis updates~\citep{metropolis1953,hastings1970} are known to slow
dramatically near criticality, while cluster methods such as Swendsen-Wang and
Wolff mitigate this problem~\citep{swendsen1987,wolff1989,sokal1997,
newman1999,binder2010}.  General reviews of machine learning in physics provide
the broader context for why careful benchmark design matters here
too~\citep{carleo2019,mehta2019}.  Relative to that prior work, the main
contribution of the present paper is methodological: a matched-budget benchmark
that keeps same-scale evidence, exploratory transfer, and sampler diagnostics
explicitly separated, and that shows how easily broad RG language can outrun
what a local majority-vote task actually supports.

% ===== SECTION 8: CONCLUSION =====
\section{Conclusion}
\label{sec:conclusion}

This revision supports a scoped negative result.  In the unified 10-seed
same-scale benchmark, the linear baseline has lower mean test MSE than the MLP
at the critical point: at $(L,\beta)=(16,\beta_c)$ the MLP reaches
$0.6249 \pm 0.6456$, while Linear reaches $0.1736 \pm 0.0862$.  The mean
ordering favors the linear model, but the strength of the difference depends on
the test because the MLP variance is large (Welch $p=0.055$, Mann-Whitney
$p=0.026$, Cohen's $d=-0.98$).  Across the benchmark as a whole, no setting
shows a robust MLP advantage over the linear baseline.  The supported claim is
therefore specific: on majority-vote block-spin tasks dominated by a local
linear-threshold structure, standard fully connected depth does not provide an
RG-like advantage over simpler linear projections.

The exploratory whole-lattice tiled transfer and the sampler diagnostics mainly
clarify methodology rather than add parallel primary evidence.  At $\beta_c$,
the tiled transfer protocol does not exhibit monotone degradation with scale
distance: $8\to16$ and even $4\to16$ remain competitive with the same-budget
$16\to16$ baseline, and linear models stay competitive throughout.  At the same
time, the critical-point sampler diagnostics show substantial autocorrelation,
with $\tau_{\mathrm{int}}(E)=18.8 \pm 7.3$ and
$\tau_{\mathrm{int}}(|M|)=37.0 \pm 15.1$ sweeps at $(16,\beta_c)$, implying
effective sample sizes of only about $80$ and $41$ for the diagnostic run.  We
therefore treat R2 as unresolved rather than falsified, and we do not promote
the appendix-level turbulence note or the older random-label pilot to primary
evidence.

The durable contribution of the paper is thus methodological.  A convincing
positive claim of RG-like computation should align matched training budgets,
true multi-scale chains (for example $4\to8\to16$ together with direct
$4\to16$ and $8\to16$ transfer), autocorrelation-aware sampling, and a regime
in which competitive linear baselines actually fail.  All main tables in this
revision are sourced from the separated raw outputs in
\path{outputs/rg_bench/cross_scale/},
\path{outputs/rg_bench/whole_lattice_transfer/}, and
\path{outputs/rg_bench/mixing_analysis/}, with the legacy pilot outputs kept
appendix-only.  On the present local majority-vote block-spin benchmark, that
combination of evidence does not support a claim that standard MLPs learn
RG-like advantages.

\appendix

\section{Sampling Diagnostics}
\label{app:mixing}

To quantify the sampling caveat rather than leave it qualitative, we ran a
dedicated time-series diagnostic using 3 seeds, 3000 post-burn-in sweeps, and
the observables $E$ and $|M|$.  Table~\ref{tab:mixing} reports the resulting
integrated autocorrelation times and mean acceptance rates.

\begin{table}[ht]
    \centering
    \caption{Metropolis-Hastings mixing diagnostics from the
        \texttt{mixing\_analysis/summary.json} output.}
    \begin{tabular}{lccc}
        \toprule
        \textbf{Setting} & \textbf{$\tau_{\mathrm{int}}(E)$} & \textbf{$\tau_{\mathrm{int}}(|M|)$} & \textbf{Acceptance} \\
        \midrule
        $L=8$, $\beta=0.30$   & $1.7 \pm 0.4$  & $2.7 \pm 0.4$  & $0.520 \pm 0.004$ \\
        $L=8$, $\beta_c$      & $6.9 \pm 1.9$  & $9.6 \pm 3.0$  & $0.171 \pm 0.005$ \\
        $L=8$, $\beta=0.60$   & $1.8 \pm 0.3$  & $1.9 \pm 0.3$  & $0.0250 \pm 0.0006$ \\
        $L=16$, $\beta=0.30$  & $1.5 \pm 0.2$  & $3.1 \pm 0.2$  & $0.522 \pm 0.001$ \\
        $L=16$, $\beta_c$     & $18.8 \pm 7.3$ & $37.0 \pm 15.1$ & $0.184 \pm 0.006$ \\
        $L=16$, $\beta=0.60$  & $1.9 \pm 0.2$  & $2.3 \pm 0.3$  & $0.0259 \pm 0.0008$ \\
        \bottomrule
    \end{tabular}
    \label{tab:mixing}
\end{table}

The critical-point autocorrelation times are an order of magnitude larger than
their off-critical counterparts, especially for $|M|$ at $L=16$.  This
confirms that the benchmark should be read as a finite-budget computational
study, not as an effectively independent-sample Monte Carlo analysis.  With
3000 post-burn-in sweeps, the corresponding effective sample sizes are
approximately $3000/(2\tau_{\mathrm{int}})\approx 218$ for $E$ and $156$ for
$|M|$ at $(L,\beta)=(8,\beta_c)$, and approximately $80$ and $41$ at
$(16,\beta_c)$.

\section{Whole-Lattice Transfer Across Temperatures}
\label{app:whole_lattice}

Table~\ref{tab:whole_lattice_all_beta} reports the reduced-budget tiled
whole-lattice transfer for all three temperatures.  The non-monotone transfer
behavior is therefore not confined to $\beta_c$, although its interpretation
remains exploratory because the full transfer experiment is only available at
$n=3$ seeds and a smaller training budget than the main benchmark.

\begin{table}[ht]
    \centering
    \caption{Exploratory whole-lattice tiled transfer across temperatures from
        \texttt{whole\_lattice\_transfer/summary.json}.  All entries use
        $n=3$ seeds, $N_{\mathrm{train}}=200$, $N_{\mathrm{test}}=100$, and
        40 epochs.}
    \resizebox{\textwidth}{!}{%
    \begin{tabular}{llccc}
        \toprule
        \textbf{$\beta$} & \textbf{Source $\to$ Target} & \textbf{MLP} & \textbf{Linear} & \textbf{Comment} \\
        \midrule
        $0.30$ & $4\to8$   & $0.088 \pm 0.044$ & $0.562 \pm 0.104$ & MLP transfer better than same-$L=8$ \\
        $0.30$ & $4\to16$  & $0.084 \pm 0.037$ & $0.545 \pm 0.081$ & MLP strongly below same-$L=16$ \\
        $0.30$ & $8\to16$  & $0.344 \pm 0.015$ & $0.509 \pm 0.013$ & Both beat same-$L=16$ \\
        \midrule
        $\beta_c$ & $4\to8$   & $0.095 \pm 0.057$ & $0.168 \pm 0.088$ & Non-monotone vs.\ source scale \\
        $\beta_c$ & $4\to16$  & $0.141 \pm 0.107$ & $0.222 \pm 0.113$ & Better than same-$L=16$ for both \\
        $\beta_c$ & $8\to16$  & $0.293 \pm 0.190$ & $0.244 \pm 0.129$ & Competitive with same-$L=16$ \\
        \midrule
        $0.60$ & $4\to8$   & $0.019 \pm 0.009$ & $0.040 \pm 0.007$ & MLP near same-$L=8$; Linear degrades \\
        $0.60$ & $4\to16$  & $0.016 \pm 0.011$ & $0.046 \pm 0.012$ & Strong model asymmetry \\
        $0.60$ & $8\to16$  & $0.016 \pm 0.013$ & $0.017 \pm 0.012$ & Linear near same-$L=16$ \\
        \bottomrule
    \end{tabular}}
    \label{tab:whole_lattice_all_beta}
\end{table}

\section{Legacy Pilot Sanity Checks}
\label{app:pilot}

Earlier fixed-split pilot experiments (random-label controls, width sweeps, and
the first ablation figures) were generated under a separate 3-seed protocol and
now live in \path{outputs/rg_bench/pilot_ablation/}.  They remain useful as
sanity checks, but they are excluded from the main quantitative tables because
their data generation, seeding, and training budgets do not match the unified
same-scale benchmark.

\section{Exploratory Turbulence Note}
\label{app:turbulence}

The turbulence closure experiment is retained only as a compressed appendix
note.  Using D2Q9 Lattice Boltzmann data at Re$=200$, the exploratory spectrum
fit gives $\alpha = -1.30 \pm 0.05$ over $k \in [2,20]$, compared with the
Kolmogorov target $-5/3 \approx -1.67$.  The gap is substantial, the inertial
range is too short at this Reynolds number, and the experiment therefore does
not support R3.  It remains a feasibility check for future higher-Re studies.

% ===== REFERENCES =====
\clearpage
\bibliography{references}

\end{document}
